\pagestyle{empty}
\tight
\begin{tblr}{width=\textwidth,colspec={|Q[c,m,1.9cm]|X[l,m]|},hlines,vlines,hspan=minimal,row{1}={bg=headblue},rowsep=1pt,colsep=3pt}
\SetCell{c,m}\hfil\logo\hfil & \textbf{POLITEKNIK NEGERI MALANG}\par\textbf{JURUSAN TEKNOLOGI INFORMASI}\par\textbf{PROGRAM STUDI: D4 TEKNIK INFORMATIKA} \\
\SetCell[c=2]{c} \textbf{RENCANA PEMBELAJARAN SEMESTER (RPS)} & \\
\end{tblr}
\begin{longtblr}{width=\textwidth,colspec={|Q[l,m,1.9cm]|X[0.85,l,m]|X[1.45,l,m]|X[0.75,l,m]|X[1.15,l,m]|},hlines,vlines,hspan=minimal,rowsep=1pt,colsep=2pt,row{1,3}={bg=headgray,font=\bfseries}}
MATA KULIAH & KODE & BOBOT (SKS)/jam & SEMESTER & TGL. PENYUSUNAN \\
Critical Thinking dan Problem Solving & RTI251003 & 2 SKS/4 jam & 1 & 6 Agustus 2025 \\
OTORISASI & Dosen Pengembang RPS & Koordinator RMK & \SetCell[c=2]{l} Ka PRODI &  \\
 & Vivi Nur Wijayaningrum, S.Kom., M.Kom. & Vivi Nur Wijayaningrum, S.Kom., M.Kom. & \SetCell[c=2]{l}{\vspace{8mm}\par Dr. Ely Setyo Astuti, S.T., M.T.} &  \\
\SetCell[r=8]{l} Capaian Pembelajaran (CP) & \SetCell[c=4]{l,bg=headgray,font=\bfseries} Capaian Pembelajaran Lulusan Program Studi (CPL-Prodi) &  &  &  \\
 & CPL04 & \SetCell[c=3]{l} Mampu menyusun deskripsi saintifik hasil kajian, pengembangan, atau implementasi teknologi informasi dan komunikasi dalam bentuk skripsi, laporan tugas akhir, atau artikel ilmiah. &  &  \\
 & \SetCell[c=4]{l,bg=headgray,font=\bfseries} Capaian Pembelajaran mata kuliah (CPL-MK) &  &  &  \\
 & CPMK0401 & \SetCell[c=3]{l} Mahasiswa mampu menganalisis permasalahan kompleks dan merumuskan solusi rasional berbasis data yang dituangkan dalam deskripsi ilmiah. &  &  \\
 & \SetCell[c=4]{l,bg=headgray,font=\bfseries} Kemampuan akhir tiap tahapan belajar (Sub-CPMK) &  &  &  \\
 & SCPMK0401-00301 & \SetCell[c=3]{l} Mahasiswa mampu mengidentifikasi, menganalisis, dan mengevaluasi argumen dalam berbagai konteks diskusi atau materi tertulis, dengan membedakan antara fakta dan opini serta mengenali bias dan kesalahan logika yang umum. &  &  \\
 & SCPMK0401-00302 & \SetCell[c=3]{l} Mahasiswa mampu menerapkan berbagai teknik dan alat pemecahan masalah, termasuk pemikiran lateral, brainstorming, dan analisis root cause, untuk menemukan solusi inovatif dan efektif terhadap masalah yang kompleks. &  &  \\
 & SCPMK0401-00303 & \SetCell[c=3]{l} Mahasiswa mampu mengembangkan solusi yang logis dan berdasarkan bukti untuk masalah yang diberikan, dan secara efektif membela solusi tersebut melalui argumentasi yang kohesif dan berbasis data dalam bentuk deskripsi ilmiah. &  &  \\
\SetCell[r=2]{l} Penilaian & \SetCell[c=4]{l} *Nilai CPMK didapat dari agregasi nilai SUB-CPMK yang dibebankan &  &  &  \\
 & \SetCell[c=4]{l}{\tiny\begin{tblr}{width=\linewidth,colspec={|Q[c,m,0.1\linewidth]|Q[c,m,0.16\linewidth]|X[l,m]|X[l,m]|X[l,m]|X[l,m]|X[l,m]|Q[c,m,0.08\linewidth]|},hlines,vlines,hspan=minimal,rowsep=1pt,colsep=0.7pt,row{1,2}={bg=headgray,font=\bfseries}}
Kode CPMK & Kode SCPMK & \SetCell[c=5]{c} Bobot per Bentuk Penilaian & & & & & Total Bobot \\
 & & Tugas & Kuis & UTS & PBL & UAS & \\
\SetCell[r=1]{c} \SetCell[r=3]{c} CPMK0401 & SCPMK0401-00301 & 10\%: analisis argumen, fakta vs opini, bias, dan kesalahan logika & 15\%: konsep berpikir kritis pekan 1--4 & 5\%: evaluasi argumen dan validitas penalaran & & & 30\% \\
\SetCell[r=2]{c}  & SCPMK0401-00302 & 8\%: teknik dan alat pemecahan masalah & & 5\%: penerapan teknik problem solving pada kasus & & & 13\% \\
 & SCPMK0401-00303 & 2\%: perumusan solusi ilmiah berbasis data & & & 20\%: draf dokumen ilmiah solusi permasalahan & 35\%: presentasi akhir solusi berbasis bukti & 57\% \\
\SetCell[c=7]{r}\textbf{Total Per Penilaian} & & & & & & & \textbf{100\%} \\
\end{tblr}} &  &  &  \\
Diskripsi Singkat Mata Kuliah & \SetCell[c=4]{l} Pada mata kuliah ini mahasiswa akan belajar mengenai konsep berpikir kritis dan bagaimana penerapannya dalam menanggapi informasi dalam kehidupan sehari-hari. Selain itu pada mata kuliah ini akan diajarkan mengenai konsep dan teknik-teknik dalam menyelesaikan kasus/permasalahan baik dalam lingkup ujian/tes, maupun pada masalah-masalah yang dijumpai dalam kehidupan sehari-hari. &  &  &  \\
Bahan Kajian / Materi Pembelajaran & \SetCell[c=4]{l}\begin{enumerate}\item Berpikir dan Menalar\item Pondasi Berpikir Kritis\item Kemampuan Dasar Pemecahan Masalah\item Applied Critical Thinking\item Kemampuan Pemecahan Masalah Tingkat Lanjut\item Teknik-teknik Lanjutan untuk Pemecahan Masalah\item Penalaran Kritis\end{enumerate} &  &  &  \\
\SetCell[r=2]{l} Pustaka & \SetCell[c=4]{l} \textbf{Utama:}\begin{enumerate}\item Thinking Skills: Critical Thinking and Problem Solving. Second Edition.\end{enumerate} &  &  &  \\
 & \SetCell[c=4]{l} \textbf{Pendukung:}\begin{enumerate}\item Critical Thinking Skills For Dummies.\end{enumerate} &  &  &  \\
Media Pembelajaran & \SetCell[c=2]{l} \textbf{Software:} Microsoft Office Word, Microsoft Office PowerPoint. &  & \SetCell[c=2]{l} \textbf{Hardware:} Laptop/PC, LCD Proyektor. &  \\
Nama Dosen Pengampu & \SetCell[c=4]{l}\begin{enumerate}\item Vivi Nur Wijayaningrum, S.Kom., M.Kom.\item Elok Nur Hamdana, S.T., M.T.\item Rokhimatul Wakhidah, S.Pd., M.T.\item Mungki Astiningrum, S.T., M.Kom.\item Bagas Satya Dian Nugraha, S.T., M.T.\end{enumerate} &  &  &  \\
Matakuliah Syarat & \SetCell[c=4]{l} - &  &  &  \\
\end{longtblr}
\begin{longtblr}{width=\textwidth,colspec={|Q[c,m,0.06\textwidth]|X[1.35,l,m]|X[1.08,l,m]|X[1.16,l,m]|X[0.7,l,m]|X[1.24,l,m]|X[1.06,l,m]|X[0.98,l,m]|Q[c,m,0.05\textwidth]|},hlines,vlines,hspan=minimal,rowhead=2,row{1,2}={bg=headgreen,font=\bfseries},rowsep=1pt,colsep=1.5pt}
Mgg.\par Ke & Kemampuan Akhir Yang Direncanakan (Sub-CP-MK) & Bahan kajian (Materi Pembelajaran) & Bentuk dan Metode Pembelajaran & Estimasi Waktu & Pengalaman Belajar Mahasiswa & Kriteria \& Bentuk Penilaian & Indikator Penilaian & Bobot Penilaian (\%) \\
(1) & (2) & (3) & (4) & (5) & (6) & (7) & (8) & (9) \\
1 & SCPMK0401-00301 & Konsep berpikir; pengertian Critical Thinking \& Problem Solving; karakteristik pemikir kritis; peran CTPS dalam kehidupan profesional. & Bentuk: Daring atau Luring. Metode: Contextual Teaching and Learning (CTL), Problem Based Learning. Penugasan: Tugas 1 contoh penerapan CTPS di kehidupan nyata. & 2 x 2 x 50' tatap muka dan tugas terstruktur. & Mahasiswa dapat memahami konsep Critical Thinking \& Problem Solving dan mengaitkannya dengan pengalaman pribadi/nyata. & Kriteria: ketepatan dan penguasaan. Bentuk: keaktifan diskusi kelompok dan ketepatan jawaban tugas. & Memahami definisi dan manfaat CTPS; mampu memberikan contoh penerapan CTPS. & 2 \\
2 & SCPMK0401-00301 & Fakta \& opini; klaim; struktur argumen dasar; valid/invalid reasoning. & Bentuk: Daring atau Luring. Metode: CTL, Problem Based Learning. Penugasan: Tugas 2 klasifikasi pernyataan sebagai fakta, opini, atau klaim. & 2 x 2 x 50' tatap muka dan tugas terstruktur. & Mahasiswa dapat menganalisis pernyataan sederhana untuk mengidentifikasi klaim dan fakta. & Kriteria: ketepatan dan penguasaan. Bentuk: keaktifan diskusi kelompok dan ketepatan jawaban tugas. & Memahami perbedaan fakta dan opini; mengidentifikasi klaim dan unsur argumen. & 2 \\
3 & SCPMK0401-00301 & Definisi masalah; identifikasi informasi; strategi pemecahan masalah; studi kasus sederhana. & Bentuk: Daring atau Luring. Metode: CTL, Problem Based Learning. Penugasan: Tugas 3 merumuskan masalah dari studi kasus. & 2 x 2 x 50' tatap muka dan tugas terstruktur. & Mahasiswa dapat mengidentifikasi masalah dan merumuskan tujuan pemecahan. & Kriteria: ketepatan dan penguasaan. Bentuk: keaktifan diskusi kelompok dan ketepatan jawaban tugas. & Menjelaskan rumusan masalah; mengidentifikasi relevansi informasi yang dipilih. & 2 \\
4 & SCPMK0401-00301 & Analisis argumen; logical evaluation; identifikasi bias; kesalahan penalaran sederhana. & Bentuk: Daring atau Luring. Metode: CTL, Problem Based Learning. Penugasan: Tugas 4 mengevaluasi argumen nyata dari studi kasus. & 2 x 2 x 50' tatap muka dan tugas terstruktur. & Mahasiswa dapat menganalisis struktur argumen. & Kriteria: ketepatan dan penguasaan. Bentuk: keaktifan diskusi kelompok dan ketepatan jawaban tugas. & Mengidentifikasi premis dan menjelaskan kesimpulan. & 2 \\
5 & SCPMK0401-00301 & Kuis materi pekan 1--4. & Ujian teori pilihan ganda, sifat closed book. & 2 x 2 x 50' tatap muka. & Mampu mengerjakan soal ujian dengan materi pekan 1--4 secara mandiri. & Kriteria: ketepatan dan penguasaan. Bentuk: ketepatan jawaban ujian. & Ketepatan jawaban dengan kunci jawaban. & 15 \\
6 & SCPMK0401-00302 & Penggunaan imajinasi; integrasi beberapa teknik; identifikasi informasi relevan; inferensi dalam problem solving. & Bentuk: Daring atau Luring. Metode: CTL, Project Based Learning. Penugasan: Tugas 5 menganalisis kasus kompleks. & 2 x 2 x 50' tatap muka dan tugas terstruktur. & Mahasiswa dapat memecahkan kasus secara bertahap. & Kriteria: ketepatan dan penguasaan. Bentuk: keaktifan diskusi kelompok dan ketepatan jawaban tugas. & Mengidentifikasi relevansi informasi yang dipilih; memahami logika inferensi. & 2 \\
7 & SCPMK0401-00302 & Mathematical reasoning; graphical reasoning; probability dan uncertainty; decision trees. & Bentuk: Daring atau Luring. Metode: CTL, Project Based Learning. Penugasan: Tugas 6 menyelesaikan permasalahan dengan pemilihan teknik yang tepat. & 2 x 2 x 50' tatap muka dan tugas terstruktur. & Mahasiswa dapat menerapkan teknik yang berbeda untuk menyelesaikan permasalahan. & Kriteria: ketepatan dan penguasaan. Bentuk: keaktifan diskusi kelompok dan ketepatan jawaban tugas. & Mengidentifikasi ketepatan teknik yang digunakan; menjelaskan konsistensi logika yang digunakan. & 2 \\
8 & SCPMK0401-00301, SCPMK0401-00302 & UTS materi pekan 1--7. & Ujian teori pilihan ganda, sifat closed book. & 2 x 2 x 50' tatap muka. & Mampu mengerjakan soal ujian dengan materi pekan 1--7 secara mandiri. & Kriteria: ketepatan dan penguasaan. Bentuk: ketepatan jawaban ujian. & Ketepatan jawaban dengan kunci jawaban. & 10 \\
9 & SCPMK0401-00301 & Validitas logika; non-deductive reasoning; argument by analogy; statistical reasoning. & Bentuk: Daring atau Luring. Metode: CTL, Project Based Learning. Penugasan: Tugas 7 memberikan kritik argumen statistik pada studi kasus. & 2 x 2 x 50' tatap muka dan tugas terstruktur. & Mahasiswa dapat menilai validitas argumen. & Kriteria: ketepatan dan penguasaan. Bentuk: keaktifan diskusi kelompok dan ketepatan jawaban tugas. & Memahami ketepatan analisis logika yang digunakan; menjelaskan argumen secara kritis berbasis data. & 2 \\
10 & SCPMK0401-00302 & Identifikasi masalah; observasi dan wawancara; 5-Why's; fishbone; problem tree. & Bentuk: Daring atau Luring. Metode: CTL, Project Based Learning. Penugasan: Tugas 8 menganalisis akar masalah dari kasus kehidupan sehari-hari. & 2 x 2 x 50' tatap muka dan tugas terstruktur. & Mahasiswa dapat mengolah data nyata atau simulasi. & Kriteria: ketepatan dan penguasaan. Bentuk: keaktifan diskusi kelompok dan ketepatan jawaban tugas. & Menjelaskan akar masalah dari sebuah kasus; menjelaskan konsistensi logika sebab-akibat. & 2 \\
11 & SCPMK0401-00302 & Brainstorming; evaluasi solusi; kreativitas terstruktur. & Bentuk: Daring atau Luring. Metode: CTL, Project Based Learning. Penugasan: Tugas 9 menyelesaikan permasalahan melalui brainstorming secara terstruktur. & 2 x 2 x 50' tatap muka dan tugas terstruktur. & Mahasiswa dapat menyusun solusi dan alternatifnya untuk menyelesaikan permasalahan. & Kriteria: ketepatan dan penguasaan. Bentuk: keaktifan diskusi kelompok dan ketepatan jawaban tugas. & Mengidentifikasi relevansi solusi berdasarkan permasalahan; menjelaskan kreativitas solusi secara terukur. & 2 \\
12 & SCPMK0401-00303 & Penyusunan solusi ilmiah; integrasi data dan argumen; validitas solusi. & Bentuk: Daring atau Luring. Metode: CTL, Project Based Learning. Penugasan: Tugas 10 menulis solusi permasalahan secara ilmiah dalam bentuk draf dokumen ilmiah. & 2 x 2 x 50' tatap muka dan tugas terstruktur. & Mahasiswa dapat menghubungkan keterkaitan antara data dan solusi. & Kriteria: ketepatan dan penguasaan. Bentuk: keaktifan diskusi kelompok dan ketepatan jawaban tugas. & Menjelaskan kekuatan argumen berdasarkan data yang diperoleh; menjelaskan kecocokan solusi untuk menyelesaikan permasalahan. & 2 \\
13 & SCPMK0401-00303 & Struktur proposal; kerangka argumentasi ilmiah. & Bentuk: Daring atau Luring. Metode: CTL, Project Based Learning (penyusunan draf dokumen ilmiah). & 2 x 2 x 50' tatap muka. & Mahasiswa dapat menyusun kerangka draf dokumen ilmiah untuk solusi permasalahan. & Kriteria: ketepatan dan penguasaan. Bentuk: keaktifan diskusi kelompok, progres draf dokumen ilmiah. & Menjelaskan solusi dari permasalahan secara logis dan terstruktur. & 5 \\
14 & SCPMK0401-00303 & Argumentasi ilmiah; integrasi data pendukung. & Bentuk: Daring atau Luring. Metode: CTL, Project Based Learning (penyusunan draf dokumen ilmiah). & 2 x 2 x 50' tatap muka. & Mahasiswa dapat menuangkan argumentasi berbasis data ke dalam draf dokumen ilmiah. & Kriteria: ketepatan dan penguasaan. Bentuk: keaktifan diskusi kelompok, progres draf dokumen ilmiah. & Menjelaskan solusi dari permasalahan secara logis dan terstruktur. & 5 \\
15 & SCPMK0401-00303 & Keterpaduan logika; penyelarasan rumusan masalah, data, dan solusi. & Bentuk: Daring atau Luring. Metode: CTL, Project Based Learning (penyusunan draf dokumen ilmiah). & 2 x 2 x 50' tatap muka. & Mahasiswa dapat menjelaskan solusi permasalahan dalam bentuk draf dokumen ilmiah yang utuh. & Kriteria: ketepatan dan penguasaan. Bentuk: keaktifan diskusi kelompok, progres draf dokumen ilmiah. & Menjelaskan solusi dari permasalahan secara logis dan terstruktur. & 5 \\
16 & SCPMK0401-00303 & Evaluasi kritis; feedback; penyempurnaan solusi. & Bentuk: Daring atau Luring. Metode: CTL, Project Based Learning (finalisasi draf dokumen ilmiah). & 2 x 2 x 50' tatap muka dan tugas terstruktur. & Mahasiswa dapat menyusun draf dokumen ilmiah secara mendalam berdasarkan analisis kritis. & Kriteria: ketepatan dan penguasaan. Bentuk: keaktifan diskusi kelompok, kualitas draf akhir dokumen ilmiah. & Menjelaskan analisis solusi permasalahan secara mendalam berdasarkan data dan fakta. & 5 \\
17 & SCPMK0401-00303 & UAS presentasi akhir. & Ujian presentasi kelompok. & 2 x 2 x 50' tatap muka. & Mahasiswa dapat mempresentasikan hasil analisis permasalahan dan solusi yang dikembangkan secara sistematis, serta menjawab pertanyaan dan memberikan klarifikasi berbasis bukti. & Kriteria: kejelasan struktur penyajian solusi, kualitas argumentasi dan keterkaitan dengan data. Bentuk: keaktifan dalam presentasi lisan dan proses tanya jawab. & Menjelaskan secara runtut latar belakang, perumusan masalah, dan solusi yang diajukan; memahami ide dan gagasan yang diusulkan untuk menyelesaikan permasalahan. & 35 \\
\end{longtblr}
